\chapter{Ring of Integers of a Quadratic Field (Human-Oriented)}

Let $d$ be a squarefree integer with $d \neq 1$. Write
$K = \Q(\sqrt{d})$.
The purpose of this chapter is to explain the classical proof that the ring of
integers of $K$ is either $\Z[\sqrt{d}]$ or
$\Z\left[\frac{1 + \sqrt{d}}{2}\right]$.

\begin{human_lemma}\label{human_orders}
  The two natural order candidates in $K$ are
  \[
  \Z[\sqrt{d}]
  \qquad\text{and}\qquad
  \Z\left[\frac{1+\sqrt{d}}{2}\right].
  \]
\end{human_lemma}
\begin{proof}
  The first is always generated by $\sqrt{d}$ over $\Z$.
  The second is relevant exactly when the half-integer
  $\frac{1+\sqrt{d}}{2}$ is integral.
\end{proof}

\begin{human_lemma}\label{human_trace_norm}
  If $x = a + b\sqrt{d} \in K$, then
  \[
  \mathrm{Tr}_{K/\Q}(x) = 2a
  \qquad\text{and}\qquad
  N_{K/\Q}(x) = a^2 - db^2.
  \]
\end{human_lemma}
\begin{proof}
  The conjugate of $x$ is $a - b\sqrt{d}$, so the trace and norm are the sum
  and product of $x$ and its conjugate.
\end{proof}

\begin{human_lemma}\label{human_integral_trace_norm}
  If $x \in \OO_K$, then both $\mathrm{Tr}_{K/\Q}(x)$ and $N_{K/\Q}(x)$ are integers.
\end{human_lemma}
\begin{proof}
  Trace and norm of an algebraic integer are integers.
\end{proof}

\begin{human_lemma}\label{human_half_int}
  If $x \in \OO_K$, then
  \[
  x = \frac{a' + b' \sqrt{d}}{2}
  \]
  for some integers $a', b'$ satisfying
  \[
  4 \mid (a'^2 - d b'^2).
  \]
\end{human_lemma}
\begin{proof}
  By the previous lemma, $2a$ is an integer and the norm is an integer.
  Setting $a' = 2a$ and $b' = 2b$ yields the displayed normal form and the
  divisibility condition
  \[
  a'^2 - d b'^2 = 4(a^2 - db^2).
  \]
\end{proof}

\begin{human_lemma}\label{human_mod4_nonone}
  If $d \not\equiv 1 \pmod 4$ and
  $4 \mid (a'^2 - d b'^2)$, then both $a'$ and $b'$ are even.
\end{human_lemma}
\begin{proof}
  A square modulo $4$ is either $0$ or $1$.
  When $d \not\equiv 1 \pmod 4$, the divisibility condition rules out the odd
  cases, so both numerators are even.
\end{proof}

\begin{human_theorem}\label{human_nonone_branch}
  If $d \not\equiv 1 \pmod 4$, then
  \[
  \OO_K = \Z[\sqrt{d}].
  \]
\end{human_theorem}
\begin{proof}
  The inclusion $\Z[\sqrt{d}] \subseteq \OO_K$ is immediate because
  $\sqrt{d}$ is integral.
  Conversely, if $x \in \OO_K$, Lemma \ref{human_half_int} gives
  $x = \frac{a' + b'\sqrt{d}}{2}$ with
  $4 \mid (a'^2 - d b'^2)$.
  Lemma \ref{human_mod4_nonone} forces $a', b'$ to be even, so $x \in \Z[\sqrt{d}]$.
\end{proof}

\begin{human_lemma}\label{human_mod4_one}
  Suppose $d \equiv 1 \pmod 4$.
  Then
  \[
  4 \mid (a'^2 - d b'^2)
  \]
  is equivalent to $a'$ and $b'$ having the same parity.
\end{human_lemma}
\begin{proof}
  Again use that squares modulo $4$ are $0$ or $1$.
  When $d \equiv 1 \pmod 4$, the divisibility condition is satisfied exactly
  when the two numerators are both even or both odd.
\end{proof}

\begin{human_lemma}\label{human_omega_integral}
  If $d \equiv 1 \pmod 4$, then
  \[
  \omega := \frac{1 + \sqrt{d}}{2}
  \]
  is integral.
\end{human_lemma}
\begin{proof}
  Writing $d = 1 + 4k$, the element $\omega$ satisfies
  \[
  \omega^2 - \omega - k = 0,
  \]
  so it is a root of a monic polynomial in $\Z[X]$.
\end{proof}

\begin{human_theorem}\label{human_one_branch}
  If $d \equiv 1 \pmod 4$, then
  \[
  \OO_K = \Z\left[\frac{1+\sqrt{d}}{2}\right].
  \]
\end{human_theorem}
\begin{proof}
  Lemma \ref{human_omega_integral} gives the inclusion
  \[
  \Z\left[\frac{1+\sqrt{d}}{2}\right] \subseteq \OO_K.
  \]
  Conversely, let $x \in \OO_K$.
  By Lemma \ref{human_half_int}, we may write
  \[
  x = \frac{a' + b'\sqrt{d}}{2}
  \]
  with $4 \mid (a'^2 - d b'^2)$.
  Lemma \ref{human_mod4_one} says that $a'$ and $b'$ have the same parity,
  which is exactly the coordinate condition for an element of
  $\Z\left[\frac{1+\sqrt{d}}{2}\right]$.
\end{proof}
